\documentclass[12pt,a4paper]{article}

\usepackage[T2A]{fontenc}
\usepackage[utf8]{inputenc}
\usepackage[russian]{babel}
\usepackage{amsmath, amssymb}
\usepackage{geometry}
\usepackage{tikz}
\usepackage{xcolor}
\usepackage{booktabs}

\usetikzlibrary{
  arrows.meta, positioning, shapes.geometric,
  fit, backgrounds, calc,
  decorations.pathreplacing
}

\geometry{top=1.5cm, bottom=1.5cm, left=1.8cm, right=1.8cm}

% ── Colors ──
\definecolor{tokblue}  {RGB}{173,216,230}
\definecolor{hidgreen} {RGB}{180,240,180}
\definecolor{outred}   {RGB}{255,200,200}
\definecolor{warnamber}{RGB}{255,220,150}
\definecolor{normpur}  {RGB}{220,200,255}
\definecolor{projcyan} {RGB}{180,230,230}
\definecolor{lossred}  {RGB}{240,128,128}
\definecolor{textcol}  {RGB}{255,228,196}
\definecolor{simgreen} {RGB}{144,238,144}
\definecolor{darkbg}   {RGB}{26,26,46}

\pagestyle{empty}

\begin{document}

\begin{center}
{\LARGE\bfseries Архитектуры ZSL/FSL: все схемы}\\[0.3em]
{\small Приложение к учебному гайду}
\end{center}

\bigskip

% ================================================================
%  СХЕМА 1: DeViSE
% ================================================================
\section*{Схема 1. DeViSE (Frome et al., NeurIPS 2013)}

Проекция CNN-признаков в пространство Word2Vec через обучаемую линейную трансформацию.

\begin{center}
\begin{tikzpicture}[
  block/.style={rectangle, draw, rounded corners=3pt,
                minimum height=0.9cm, text centered, font=\small, thick,
                text width=2.4cm, align=center},
  arr/.style={-{Stealth[length=6pt]}, thick},
  dim/.style={font=\scriptsize, color=gray},
]

% Image path
\node[block, fill=tokblue, minimum width=2cm] (img) at (0, 0) {Изображение\\$224\times224$};
\node[block, fill=hidgreen, minimum width=2.2cm] (cnn) at (3, 0) {CNN\\(AlexNet)};
\node[dim, below=0.1cm of cnn] {4096-dim};
\node[block, fill=warnamber, minimum width=2.2cm] (proj) at (6.2, 0) {Линейная\\проекция $M$};
\node[block, fill=projcyan, minimum width=2cm] (ve) at (9.2, 0) {Визуальный\\эмбеддинг};
\node[dim, below=0.1cm of ve] {500-dim};

\draw[arr] (img) -- (cnn);
\draw[arr] (cnn) -- (proj);
\draw[arr] (proj) -- (ve);

% Text path
\node[block, fill=textcol, minimum width=2.2cm] (word) at (6.2, -2.5) {Word2Vec\\(Skip-gram)};
\node[block, fill=projcyan, minimum width=2cm] (te) at (9.2, -2.5) {Семантический\\эмбеддинг};
\node[dim, below=0.1cm of te] {500-dim};

\node[font=\small, anchor=west] at (3.5, -2.5) {Название класса <<cat>>};
\draw[arr] (4.9, -2.5) -- (word);
\draw[arr] (word) -- (te);

% Similarity
\node[block, fill=simgreen, minimum width=2cm] (sim) at (12, -1.25) {Hinge\\Rank Loss};
\draw[arr] (ve.east) -- ++(0.5,0) |- (sim.west);
\draw[arr] (te.east) -- ++(0.5,0) |- (sim.west);

% Annotation
\node[font=\scriptsize, color=gray, text width=4cm, align=center] at (12, -2.8)
  {Правильный класс ближе\\к изображению, чем неправильные\\(с отступом $\Delta = 0.1$)};

\end{tikzpicture}
\end{center}

\[
L = \sum_{j \neq y} \max\!\Big[0,\; \Delta - \tilde{t}_y \cdot M\bar{v}(x) + \tilde{t}_j \cdot M\bar{v}(x)\Big]
\]
где $\tilde{t}_y$ --- Word2Vec эмбеддинг правильного класса, $M\bar{v}(x)$ --- спроецированные CNN-признаки.


% ================================================================
%  СХЕМА 2: f-CLSWGAN
% ================================================================
\newpage
\section*{Схема 2. f-CLSWGAN (Xian et al., CVPR 2018)}

Условный WGAN генерирует синтетические CNN-признаки для unseen-классов.

\begin{center}
\begin{tikzpicture}[
  block/.style={rectangle, draw, rounded corners=3pt,
                minimum height=0.9cm, text centered, font=\small, thick,
                text width=2.5cm, align=center},
  arr/.style={-{Stealth[length=6pt]}, thick},
  dim/.style={font=\scriptsize, color=gray},
]

% Inputs to Generator
\node[block, fill=normpur, minimum width=2cm] (noise) at (0, 0) {Шум\\$z \sim \mathcal{N}(0, I)$};
\node[block, fill=textcol, minimum width=2cm] (desc) at (0, -1.8) {Описание класса\\$c(y)$};
\node[dim, below=0.05cm of desc] {атрибуты / w2v};

% Generator
\node[block, fill=hidgreen, minimum width=2.8cm, minimum height=1.5cm] (gen) at (3.8, -0.9) {\textbf{Генератор $G$}\\MLP\\(4096 hidden)};
\draw[arr] (noise.east) -- ++(0.5,0) |- ([yshift=3pt]gen.west);
\draw[arr] (desc.east) -- ++(0.5,0) |- ([yshift=-3pt]gen.west);

% Generated feature
\node[block, fill=projcyan, minimum width=2.5cm] (fake) at (7.5, -0.9) {Синтетический\\признак $\tilde{x}$};
\node[dim, below=0.1cm of fake] {2048-dim};
\draw[arr] (gen) -- (fake);

% Discriminator
\node[block, fill=outred, minimum width=2.8cm, minimum height=1.2cm] (disc) at (11, 0.3) {\textbf{Дискриминатор $D$}\\MLP\\<<реальный / сгенерированный?>>};

% Real feature
\node[block, fill=tokblue, minimum width=2.5cm] (real) at (7.5, 1.5) {Реальный признак $x$\\(ResNet-101)};
\node[dim, below=0.1cm of real] {2048-dim};

\draw[arr] (fake.east) -- ++(0.5,0) |- ([yshift=-4pt]disc.west);
\draw[arr] (real.east) -- ++(0.5,0) |- ([yshift=4pt]disc.west);

% Classification loss
\node[block, fill=lossred, minimum width=2.5cm] (cls) at (11, -2.2) {\textbf{Классификатор}\\$P(y|\tilde{x};\theta)$};
\draw[arr] (fake.east) -- ++(0.5,0) |- (cls.west);

% Loss labels
\node[font=\small\bfseries, color=red!70!black, anchor=west] at (13, 0.3) {$L_{WGAN}$};
\node[font=\small\bfseries, color=red!70!black, anchor=west] at (13, -2.2) {$L_{CLS}$};

% Total loss
\node[font=\small, text width=7cm, anchor=north west] at (0, -3.5)
  {Полная функция потерь:
   $\displaystyle\min_G \max_D\; L_{WGAN} + \beta \cdot L_{CLS}$, \quad $\beta = 0.01$, $\lambda = 10$};

\end{tikzpicture}
\end{center}


% ================================================================
%  СХЕМА 3: DAP Pipeline
% ================================================================
\newpage
\section*{Схема 3. DAP Pipeline (Lampert et al., CVPR 2009)}

Предсказание класса через промежуточные атрибуты.

\begin{center}
\begin{tikzpicture}[
  block/.style={rectangle, draw, rounded corners=3pt,
                minimum height=0.8cm, text centered, font=\small, thick,
                text width=2.2cm, align=center},
  ablock/.style={rectangle, draw, rounded corners=2pt,
                minimum height=0.55cm, text centered, font=\scriptsize, thick,
                text width=2.8cm, align=center, fill=warnamber},
  arr/.style={-{Stealth[length=5pt]}, thick},
  dim/.style={font=\scriptsize, color=gray},
]

% Input
\node[block, fill=tokblue, minimum width=2.5cm] (img) at (0, 0) {Изображение $x$};

% Attribute classifiers
\node[ablock] (a1) at (4, 1.5) {$p(a_1|x)$: полосатый?};
\node[ablock] (a2) at (4, 0.5) {$p(a_2|x)$: четвероногий?};
\node[ablock] (a3) at (4, -0.5) {$p(a_3|x)$: водоплавающий?};
\node[ablock] (am) at (4, -1.5) {$p(a_M|x)$: чёрно-белый?};

\node[font=\small, color=gray] at (4, -2.2) {$M$ бинарных классификаторов};

\draw[arr] (img.east) -- ++(0.5,0) |- (a1.west);
\draw[arr] (img.east) -- ++(0.5,0) |- (a2.west);
\draw[arr] (img.east) -- ++(0.5,0) |- (a3.west);
\draw[arr] (img.east) -- ++(0.5,0) |- (am.west);

% Predicted vector
\node[block, fill=projcyan, minimum width=2.5cm, minimum height=1.5cm] (pred) at (8, 0) {Предсказанный\\вектор\\атрибутов};

\draw[arr] (a1.east) -- ++(0.3,0) |- ([yshift=4pt]pred.west);
\draw[arr] (a2.east) -- ++(0.3,0) |- ([yshift=1pt]pred.west);
\draw[arr] (a3.east) -- ++(0.3,0) |- ([yshift=-2pt]pred.west);
\draw[arr] (am.east) -- ++(0.3,0) |- ([yshift=-5pt]pred.west);

% Comparison
\node[block, fill=simgreen, minimum width=2.5cm, minimum height=1.2cm] (cmp) at (11.5, 0) {Сравнение\\с $a^y$ каждого\\unseen-класса};

\draw[arr] (pred) -- (cmp);

% Output
\node[block, fill=outred, minimum width=2cm, thick] (out) at (14.5, 0) {$\hat{y}$\\argmax};

\draw[arr] (cmp) -- (out);

% Class descriptions
\node[font=\scriptsize, text width=3.5cm, anchor=north, color=gray] at (11.5, -1.5)
  {$a^{\text{зебра}} = (1, 1, 0, 1, \ldots)$\\
   $a^{\text{пингвин}} = (0, 0, 1, 1, \ldots)$\\
   $a^{\text{тигр}} = (1, 1, 0, 0, \ldots)$};

\end{tikzpicture}
\end{center}

\[
\hat{y} = \arg\max_{y \in Y_{unseen}} \sum_{m=1}^{M} \Big[ a_m^y \cdot \log p(a_m|x) + (1 - a_m^y) \cdot \log(1 - p(a_m|x)) \Big]
\]


% ================================================================
%  СХЕМА 5: OWL-ViT
% ================================================================
\newpage
\section*{Схема 5. OWL-ViT (Minderer et al., ECCV 2022)}

Адаптация CLIP для zero-shot object detection: каждый patch-токен = потенциальный объект.

\begin{center}
\begin{tikzpicture}[
  block/.style={rectangle, draw, rounded corners=3pt,
                minimum height=0.8cm, text centered, font=\small, thick,
                text width=2.5cm, align=center},
  arr/.style={-{Stealth[length=5pt]}, thick},
  dim/.style={font=\scriptsize, color=gray},
]

% Image input
\node[block, fill=tokblue, minimum width=2.5cm] (img) at (0, 0) {Изображение\\$768\times768$};

% ViT
\node[block, fill=hidgreen, minimum width=3cm, minimum height=1.5cm] (vit) at (3.5, 0) {\textbf{ViT Encoder}\\(CLIP backbone)\\Без final pooling!};
\draw[arr] (img) -- (vit);

% Patch tokens
\node[block, fill=projcyan, minimum width=2.8cm, minimum height=1cm] (patches) at (7.2, 0) {$N \times N$ patch\\токенов\\(576 шт.)};
\node[dim, below=0.1cm of patches] {каждый 768-dim};
\draw[arr] (vit) -- (patches);

% Classification head
\node[block, fill=warnamber, minimum width=2.8cm] (clshead) at (11, 1) {\textbf{Classification Head}\\Linear 768$\to$512};
\draw[arr] (patches.east) -- ++(0.3,0) |- (clshead.west);

% Box head
\node[block, fill=outred, minimum width=2.8cm] (boxhead) at (11, -1) {\textbf{Box Head}\\MLP $\to$ $[c_x, c_y, w, h]$};
\draw[arr] (patches.east) -- ++(0.3,0) |- (boxhead.west);

% Text queries
\node[block, fill=textcol, minimum width=2.8cm] (text) at (11, 3) {Text Encoder\\<<a cat>>, <<a dog>>};
\node[dim, below=0.05cm of text] {текстовые эмбеддинги};

% Cosine sim
\node[block, fill=simgreen, minimum width=1.5cm] (cos) at (14, 1) {$\cos$};
\draw[arr] (clshead) -- (cos);
\draw[arr] (text.east) -- ++(0.5,0) |- (cos.north);

% NMS
\node[block, fill=lossred, minimum width=2cm, thick] (nms) at (14, -1) {NMS\\$\to$ boxes};
\draw[arr] (cos.south) -- (nms.north);
\draw[arr] (boxhead) -- (nms);

% Key difference
\node[font=\small, text width=8cm, anchor=north west] at (0, -2.5)
  {\textbf{Ключевое отличие от CLIP:} CLIP берёт [CLS]-токен $\to$ один вектор на всё изображение.\\
   OWL-ViT использует \textbf{все patch-токены} $\to$ каждый патч = потенциальный объект.\\
   К каждому токену прикреплены Classification Head ($\cos$ с текстом) + Box Head ($\to$ bbox).};

\end{tikzpicture}
\end{center}


% ================================================================
%  СХЕМА 6: Siamese Networks
% ================================================================
\newpage
\section*{Схема 6. Siamese Networks (Koch et al., ICML Workshop 2015)}

Две одинаковые CNN с общими весами определяют, принадлежат ли два изображения одному классу.

\begin{center}
\begin{tikzpicture}[
  block/.style={rectangle, draw, rounded corners=3pt,
                minimum height=0.9cm, text centered, font=\small, thick,
                text width=2.5cm, align=center},
  arr/.style={-{Stealth[length=5pt]}, thick},
  dim/.style={font=\scriptsize, color=gray},
]

% Image 1 path
\node[block, fill=tokblue] (img1) at (0, 1.5) {Image$_1$};
\node[block, fill=hidgreen, minimum width=3cm] (cnn1) at (3.5, 1.5) {CNN\\(shared weights)};
\node[block, fill=projcyan] (h1) at (6.5, 1.5) {$\mathbf{h}_1$};
\node[dim, below=0.1cm of h1] {4096-dim};

\draw[arr] (img1) -- (cnn1);
\draw[arr] (cnn1) -- (h1);

% Image 2 path
\node[block, fill=tokblue] (img2) at (0, -1.5) {Image$_2$};
\node[block, fill=hidgreen, minimum width=3cm] (cnn2) at (3.5, -1.5) {CNN\\(shared weights)};
\node[block, fill=projcyan] (h2) at (6.5, -1.5) {$\mathbf{h}_2$};
\node[dim, below=0.1cm of h2] {4096-dim};

\draw[arr] (img2) -- (cnn2);
\draw[arr] (cnn2) -- (h2);

% Shared weights annotation
\draw[dashed, gray, thick] (3.5, 0.7) -- (3.5, -0.7) node[midway, right, font=\scriptsize, color=gray] {общие веса};

% Comparison
\node[block, fill=warnamber, minimum width=2.5cm] (diff) at (9.5, 0) {Взвешенная\\$|\mathbf{h}_1 - \mathbf{h}_2|$};
\draw[arr] (h1.east) -- ++(0.3,0) |- (diff.north);
\draw[arr] (h2.east) -- ++(0.3,0) |- (diff.south);

% FC + sigmoid
\node[block, fill=outred] (fc) at (12.5, 0) {FC $\to$ $\sigma$};
\draw[arr] (diff) -- (fc);

% Output
\node[font=\large\bfseries] (out) at (14.5, 0) {$p \in [0,1]$};
\draw[arr] (fc) -- (out);

\end{tikzpicture}
\end{center}

\[
p(x_1, x_2) = \sigma\!\left(\sum_j \alpha_j \cdot |h_1^{(j)} - h_2^{(j)}|\right)
\qquad
L = y \cdot \log p + (1-y) \cdot \log(1-p)
\]
$\alpha_j$ --- обучаемый вес $j$-й компоненты; $y=1$ если один класс, $y=0$ если разные.


% ================================================================
%  СХЕМА 7: Matching Networks
% ================================================================
\newpage
\section*{Схема 7. Matching Networks (Vinyals et al., NeurIPS 2016)}

Классификация через взвешенное голосование по support set с attention.

\begin{center}
\begin{tikzpicture}[
  block/.style={rectangle, draw, rounded corners=3pt,
                minimum height=0.8cm, text centered, font=\small, thick,
                text width=2.5cm, align=center},
  arr/.style={-{Stealth[length=5pt]}, thick},
  dim/.style={font=\scriptsize, color=gray},
]

% Support set
\node[block, fill=tokblue, minimum width=3cm, minimum height=1.5cm] (supp) at (0, 0)
  {Support set\\$\{(x_1,y_1), \ldots,$\\$(x_K, y_K)\}$};

% bi-LSTM
\node[block, fill=hidgreen, minimum width=2.5cm, minimum height=1cm] (lstm) at (3.5, 0)
  {\textbf{bi-LSTM}\\(FCE)};
\draw[arr] (supp) -- (lstm);

% Embeddings g
\node[block, fill=projcyan, minimum width=2.5cm] (gemb) at (6.5, 0)
  {$g(x_i, S)$\\контекстные};
\draw[arr] (lstm) -- (gemb);

% Query
\node[block, fill=textcol, minimum width=2.5cm] (query) at (0, -3)
  {Query $\hat{x}$};
\node[block, fill=hidgreen, minimum width=2.5cm] (fenc) at (3.5, -3)
  {Encoder $f$};
\node[block, fill=projcyan, minimum width=2cm] (femb) at (6.5, -3)
  {$f(\hat{x})$};

\draw[arr] (query) -- (fenc);
\draw[arr] (fenc) -- (femb);

% Cosine + softmax
\node[block, fill=simgreen, minimum width=2.5cm, minimum height=1.2cm] (attn) at (10, -1.5)
  {$\cos \to$ softmax\\$\to$ weights $a_i$};
\draw[arr] (gemb.east) -- ++(0.5,0) |- ([yshift=3pt]attn.west);
\draw[arr] (femb.east) -- ++(0.5,0) |- ([yshift=-3pt]attn.west);

% Weighted vote
\node[block, fill=outred, minimum width=2.5cm, thick] (vote) at (13.5, -1.5)
  {$\hat{y} = \sum_i a_i y_i$\\взвеш. голосование};
\draw[arr] (attn) -- (vote);

\end{tikzpicture}
\end{center}

\[
a(\hat{x}, x_i) = \frac{\exp\!\big(\cos(f(\hat{x}),\; g(x_i, S))\big)}{\sum_j \exp\!\big(\cos(f(\hat{x}),\; g(x_j, S))\big)}
\qquad
\hat{y} = \sum_{i=1}^{K} a(\hat{x}, x_i) \cdot y_i
\]

$g(x_i, S) = \overrightarrow{h}_i + \overleftarrow{h}_i + g'(x_i)$ --- Full Context Embedding через bi-LSTM.


% ================================================================
%  СХЕМА 8: Prototypical Networks
% ================================================================
\newpage
\section*{Схема 8. Prototypical Networks (Snell et al., NeurIPS 2017)}

Прототип = среднее эмбеддингов support-примеров. Классификация по ближайшему прототипу.

\begin{center}
\begin{tikzpicture}[
  block/.style={rectangle, draw, rounded corners=3pt,
                minimum height=0.8cm, text centered, font=\small, thick,
                text width=2.5cm, align=center},
  arr/.style={-{Stealth[length=5pt]}, thick},
  dim/.style={font=\scriptsize, color=gray},
]

% Support per class
\node[block, fill=tokblue, minimum width=2cm] (s1) at (0, 1.5)  {Support класс 1\\$K$ примеров};
\node[block, fill=tokblue, minimum width=2cm] (s2) at (0, 0)    {Support класс 2\\$K$ примеров};
\node[block, fill=tokblue, minimum width=2cm] (sn) at (0, -1.5) {Support класс $N$\\$K$ примеров};

% Encoder
\node[block, fill=hidgreen, minimum width=2.5cm, minimum height=3.5cm] (enc) at (3.5, 0)
  {\textbf{Encoder} $f_\phi$\\(CNN,\\4 conv блока,\\shared)};

\draw[arr] (s1) -- (enc.north west |- s1);
\draw[arr] (s2) -- (enc);
\draw[arr] (sn) -- (enc.south west |- sn);

% Mean
\node[block, fill=warnamber, minimum width=2cm] (m1) at (6.5, 1.5)  {mean $\to \mathbf{c}_1$};
\node[block, fill=warnamber, minimum width=2cm] (m2) at (6.5, 0)    {mean $\to \mathbf{c}_2$};
\node[block, fill=warnamber, minimum width=2cm] (mn) at (6.5, -1.5) {mean $\to \mathbf{c}_N$};

\draw[arr] ([yshift=8pt]enc.east) -- (m1.west);
\draw[arr] (enc.east) -- (m2.west);
\draw[arr] ([yshift=-8pt]enc.east) -- (mn.west);

% Query path
\node[block, fill=textcol, minimum width=2cm] (q) at (0, -3.5)  {Query $x^*$};
\node[block, fill=hidgreen, minimum width=2cm] (qenc) at (3.5, -3.5) {$f_\phi$\\(same)};
\draw[arr] (q) -- (qenc);

% Distance + softmax
\node[block, fill=simgreen, minimum width=3cm, minimum height=1.5cm] (dist) at (10, -1)
  {$d(f_\phi(x^*), \mathbf{c}_k)$\\$\downarrow$\\softmax $\to$ $p(y=k|x^*)$};

\draw[arr] (m1.east) -- ++(0.3,0) |- ([yshift=5pt]dist.west);
\draw[arr] (m2.east) -- ++(0.3,0) |- (dist.west);
\draw[arr] (mn.east) -- ++(0.3,0) |- ([yshift=-5pt]dist.west);
\draw[arr] (qenc.east) -- ++(3.5,0) |- ([yshift=-10pt]dist.west);

% Output
\node[block, fill=outred, thick, minimum width=2cm] (out) at (13.5, -1) {$\hat{y}$\\argmax};
\draw[arr] (dist) -- (out);

\end{tikzpicture}
\end{center}

\[
\mathbf{c}_k = \frac{1}{|S_k|} \sum_{(x_i, y_i) \in S_k} f_\phi(x_i)
\qquad
p(y\!=\!k \mid x^*) = \frac{\exp(-d(f_\phi(x^*), \mathbf{c}_k))}{\sum_{k'} \exp(-d(f_\phi(x^*), \mathbf{c}_{k'}))}
\]

$d$ --- евклидово расстояние (дивергенция Брэгмана для $\varphi(z) = \|z\|^2$).


% ================================================================
%  СХЕМА 9: MAML
% ================================================================
\newpage
\section*{Схема 9. MAML (Finn et al., ICML 2017)}

Мета-обучение инициализации через двухуровневую оптимизацию.

\begin{center}
\begin{tikzpicture}[
  block/.style={rectangle, draw, rounded corners=3pt,
                minimum height=0.8cm, text centered, font=\small, thick,
                text width=2.8cm, align=center},
  arr/.style={-{Stealth[length=5pt]}, thick},
  dim/.style={font=\scriptsize, color=gray},
]

% Meta-weights
\node[block, fill=warnamber, minimum width=3cm, minimum height=1.2cm, line width=1.5pt] (theta) at (0, 0)
  {\textbf{Мета-веса} $\theta$\\(общая инициализация)};

% Inner loops
\node[block, fill=tokblue] (t1) at (5, 2)   {Задача $T_1$\\(support set)};
\node[block, fill=tokblue] (t2) at (5, 0)   {Задача $T_2$\\(support set)};
\node[block, fill=tokblue] (t3) at (5, -2)  {Задача $T_3$\\(support set)};

\node[block, fill=hidgreen] (th1) at (9, 2)  {$\theta'_1 = \theta - \alpha\nabla L_{T_1}$};
\node[block, fill=hidgreen] (th2) at (9, 0)  {$\theta'_2 = \theta - \alpha\nabla L_{T_2}$};
\node[block, fill=hidgreen] (th3) at (9, -2) {$\theta'_3 = \theta - \alpha\nabla L_{T_3}$};

\draw[arr] (theta.east) -- ++(0.5,0) |- (t1.west);
\draw[arr] (theta.east) -- ++(0.5,0) |- (t2.west);
\draw[arr] (theta.east) -- ++(0.5,0) |- (t3.west);

\draw[arr] (t1) -- node[above, font=\tiny, color=gray]{inner loop} (th1);
\draw[arr] (t2) -- node[above, font=\tiny, color=gray]{inner loop} (th2);
\draw[arr] (t3) -- node[above, font=\tiny, color=gray]{inner loop} (th3);

% Outer loop: eval on query sets
\node[block, fill=outred, minimum width=3cm] (eval) at (13, 0) {Eval на query sets\\$\sum_i L_{T_i}(f_{\theta'_i})$};

\draw[arr] (th1.east) -- ++(0.3,0) |- ([yshift=4pt]eval.west);
\draw[arr] (th2.east) -- (eval.west);
\draw[arr] (th3.east) -- ++(0.3,0) |- ([yshift=-4pt]eval.west);

% Update theta
\draw[arr, dashed, red, thick] (eval.south) -- ++(0,-1.2) -| (theta.south)
  node[pos=0.25, below, font=\small, color=red] {outer loop: $\theta \leftarrow \theta - \beta\nabla_\theta \sum_i L_{T_i}(f_{\theta'_i})$};

% Labels
\draw[decorate, decoration={brace, amplitude=6pt}]
  (6.5, 2.6) -- (10.5, 2.6) node[midway, above=6pt, font=\small] {Inner loop ($\alpha$)};

\end{tikzpicture}
\end{center}

\begin{align*}
&\textbf{Inner:}\quad \theta'_i = \theta - \alpha \nabla_\theta L_{T_i}(f_\theta) && \text{--- адаптация к задаче}\\
&\textbf{Outer:}\quad \theta \leftarrow \theta - \beta \nabla_\theta \sum_{T_i} L_{T_i}(f_{\theta'_i}) && \text{--- мета-оптимизация}
\end{align*}


% ================================================================
%  СХЕМА 10: SetFit
% ================================================================
\newpage
\section*{Схема 10. SetFit (Tunstall et al., 2022)}

Двухэтапный few-shot: контрастивное дообучение Sentence Transformer + классификационная голова.

\begin{center}
\begin{tikzpicture}[
  block/.style={rectangle, draw, rounded corners=3pt,
                minimum height=0.8cm, text centered, font=\small, thick,
                text width=2.8cm, align=center},
  arr/.style={-{Stealth[length=5pt]}, thick},
  dim/.style={font=\scriptsize, color=gray},
]

% Stage 1
\node[font=\large\bfseries, color=blue!70!black] at (-0.5, 3.5) {Этап 1};

\node[block, fill=tokblue, minimum width=2cm] (data) at (0, 2.3) {$K$ примеров\\на класс};

\node[block, fill=warnamber, minimum width=3cm] (pairs) at (3.5, 2.3) {Генерация пар\\позитивные ($y\!=\!1$)\\негативные ($y\!=\!0$)};
\node[dim, below=0.1cm of pairs] {$R=20$ пар/класс};

\draw[arr] (data) -- (pairs);

\node[block, fill=hidgreen, minimum width=3.2cm, minimum height=1.2cm] (st) at (7.5, 2.3)
  {\textbf{Sentence}\\
   \textbf{Transformer}\\
   contrastive fine-tune};

\draw[arr] (pairs) -- (st);

% Loss
\node[block, fill=lossred, minimum width=3.5cm] (loss) at (12, 2.3)
  {$L = (\cos(f(x_i), f(x_j)) - y)^2$};

\draw[arr] (st) -- (loss);

% Stage 2
\node[font=\large\bfseries, color=green!50!black] at (-0.5, 0) {Этап 2};

\node[block, fill=projcyan, minimum width=3cm] (embs) at (3.5, 0) {Эмбеддинги из\\дообученного ST};
\draw[arr] (st.south) -- ++(0, -0.6) -| (embs.north);

\node[block, fill=outred, minimum width=2.5cm] (lr) at (7.5, 0)
  {Logistic\\Regression};
\draw[arr] (embs) -- (lr);

\node[block, fill=simgreen, minimum width=2cm, thick] (clf) at (10.5, 0) {Classifier};
\draw[arr] (lr) -- (clf);

% Inference
\node[font=\large\bfseries, color=gray] at (-0.5, -2) {Инференс};

\node[block, fill=textcol] (inp) at (3.5, -2) {Новый текст $x$};
\node[block, fill=hidgreen] (st2) at (7.5, -2) {ST (дообуч.)};
\node[block, fill=outred] (lr2) at (10.5, -2) {LR head};
\node[font=\large\bfseries] (out) at (13, -2) {$\hat{y}$};

\draw[arr] (inp) -- (st2);
\draw[arr] (st2) -- (lr2);
\draw[arr] (lr2) -- (out);

\end{tikzpicture}
\end{center}

Результат: SetFit ($\sim$355M) \textbf{превосходит} GPT-3 (175B) на RAFT: 71.3\% vs 62.7\%.


% ================================================================
%  СХЕМА 11: LoRA
% ================================================================
\newpage
\section*{Схема 11. LoRA (Hu et al., ICLR 2022)}

Низкоранговое разложение обновления весов: $W = W_0 + BA$, $r \ll \min(d,k)$.

\begin{center}
\begin{tikzpicture}[
  block/.style={rectangle, draw, rounded corners=3pt,
                minimum height=0.8cm, text centered, font=\small, thick,
                text width=2.5cm, align=center},
  arr/.style={-{Stealth[length=5pt]}, thick},
  dim/.style={font=\scriptsize, color=gray},
  frozen/.style={pattern=north east lines, pattern color=red!30},
]

% Input
\node[block, fill=tokblue, minimum width=2cm] (x) at (0, 0) {Вход $x$\\$\in \mathbb{R}^k$};

% W0 path (frozen)
\node[block, fill=outred, minimum width=3cm, minimum height=1.5cm, line width=1.5pt] (w0) at (4, 1.8)
  {\textbf{$W_0$} ($d \times k$)\\предобученные веса\\\textcolor{red}{ЗАМОРОЖЕНЫ}};
\draw[arr] (x.north) -- ++(0, 0.5) -| (w0.south);

\node[font=\small] (h0) at (7.5, 1.8) {$h_0 = W_0 x$};
\draw[arr] (w0) -- (h0);

% BA path (trainable)
\node[block, fill=hidgreen, minimum width=1.8cm] (A) at (2.5, -1.8) {$A$\\($r \times k$)};
\node[dim, below=0.1cm of A] {обучаемая};

\node[block, fill=hidgreen, minimum width=1.8cm] (B) at (5.5, -1.8) {$B$\\($d \times r$)};
\node[dim, below=0.1cm of B] {обучаемая};

\draw[arr] (x.south) -- ++(0, -0.5) -| (A.north);
\draw[arr] (A) -- (B);

\node[font=\small] (dh) at (7.5, -1.8) {$\Delta h = \frac{\alpha}{r} BAx$};
\draw[arr] (B) -- (dh);

% Sum
\node[circle, draw, fill=warnamber, minimum size=0.8cm, font=\large, thick] (sum) at (10, 0) {$+$};
\draw[arr] (h0) -- ++(0.5,0) |- (sum.north);
\draw[arr] (dh) -- ++(0.5,0) |- (sum.south);

% Output
\node[block, fill=projcyan, minimum width=2.5cm, thick] (out) at (13, 0) {Выход $h$\\$= h_0 + \Delta h$\\$\in \mathbb{R}^d$};
\draw[arr] (sum) -- (out);

% Rank annotation
\draw[decorate, decoration={brace, amplitude=6pt, mirror}]
  (1.4, -2.8) -- (6.8, -2.8)
  node[midway, below=8pt, font=\small, color=gray] {rank $r \ll \min(d, k)$, обычно $r = 4$ или $8$};

% Key properties
\node[font=\small, text width=12cm, anchor=north west] at (0, -4.2)
  {Init: $A \sim \mathcal{N}(0, \sigma^2)$, $B = 0$ $\Rightarrow$ $\Delta W = 0$ в начале (старт = предобученная модель).\\
   Инференс: $W = W_0 + BA$ вычисляется явно \textbf{один раз} $\Rightarrow$ нулевая дополнительная задержка.\\
   GPT-3 (175B): обучаемых параметров в \textbf{10\,000$\times$} меньше, GPU-памяти в \textbf{3$\times$} меньше.};

\end{tikzpicture}
\end{center}

\[
\boxed{W = W_0 + \Delta W = W_0 + BA, \quad B \in \mathbb{R}^{d \times r},\; A \in \mathbb{R}^{r \times k},\; r \ll \min(d,k)}
\]

\end{document}